% Samenvatting van het vak Vergelijkende Studie van Imperatieve Programmeertalen
\documentclass[a4paper]{article}
\usepackage[margin=2.5cm]{geometry}
\usepackage{parskip}
\usepackage[pdftex]{graphicx}
\usepackage[english]{babel}
\usepackage{amsmath}
\usepackage{float}
\usepackage{framed}
\usepackage{varwidth}
\usepackage{listings}
\usepackage{color}
\usepackage{textcomp}
\usepackage{qrcode}

\lstset{
    language=Scheme,
    frame=leftline
}
\lstMakeShortInline[columns=fixed]|

\begin{document}

\begin{titlepage}
    \newpage
    \thispagestyle{empty}
    \frenchspacing
    \hspace{-0.2cm}
    \includegraphics[height=3.4cm]{assets/sedes}
    \hspace{0.2cm}
    \rule{0.5pt}{3.4cm}
    \hspace{0.2cm}
    \begin{minipage}[b]{8cm}
        \Large{Katholieke\newline Universiteit\newline Leuven}\smallskip\newline
        \large{}\smallskip\newline
        \textbf{Department of\newline Computer Science}\smallskip
    \end{minipage}
    \hspace{\stretch{1}}
    \vspace*{3.2cm}\vfill
    \begin{center}
        \begin{minipage}[t]{\textwidth}
            \begin{center}
                \large{\rm{\textbf{\uppercase{Notes}}}}\\
                \large{\rm{Comparative Programming Languages [H0S01a]}}\\
                \large{\rm{Vergelijkende Studie van Imperatieve Programmeertalen [H04L5a]}}
            \end{center}
        \end{minipage}
    \end{center}
    \vfill\vspace*{5cm}
    \begin{minipage}[t]{0.2\textwidth}
        \begin{center}
            @GitHub\\
            \qrcode{https://github.com/FHannes/cpl-notes}\\
            {\tt \#\input{rev}}
        \end{center}
    \end{minipage}\hfill
    \begin{minipage}[t]{0.3\textwidth}
        \vfill\noindent \\[5mm]
        {\rm \textbf{Fr\'ed\'eric Hannes}}\\[2mm]
        {\rm Academic year 2015--2016}
    \end{minipage}
\end{titlepage}

%\setcounter{secnumdepth}{4}
%\setcounter{tocdepth}{4}

\newcommand{\subsubsubsection}[1]{\paragraph{#1}\mbox{}\\}

\tableofcontents

\newpage

\section{Info}

This document contains my notes for the course Comparative Programming Languages. It features code examples from the textbook used in the course, with some or a lot of explanation, depending on the complexity of the examples. I try to implement these myself first, if my implementation differs greatly from that of the textbook, I also add it to this document with an explanation of what I could've done better. After these you will also find examples, exercises and a modest Racket language reference guide.

\newpage

\section{Code}

\subsection{Chapter 1}

\subsubsection{in-S?}

The set $S$ is defined as a set containing any natural number $n$ where $n = 0$ or $n - 3 \in S$. The procedure |in-S?| returns |#t| if a given value is part of the set $S$ and |#f| if it is not.

\begin{lstlisting}
(define in-S?
  (lambda (x)
    (if (zero? x) #t
        (if (< x 3) #f
            (in-S? (- x 3))))))
\end{lstlisting}

\subsubsection{list-length}

The procedure |list-length| returns the number of elements in a given list. This functionality is available as |length| in the Racket API.

\textbf{list-length} : \textit{Listof(Sym) $\rightarrow$ Int}
\begin{lstlisting}[aboveskip=0pt]
(define list-length
  (lambda (lst)
    (if (null? lst)
        0
        (+ 1 (list-length (cdr lst))))))
\end{lstlisting}

\subsubsection{nth-element}

The procedure |nth-element| returns the element from a given list at a given position. This functionality is available as |list-ref| in the Racket API.

\textbf{nth-element} : \textit{Listof(Sym) $\times$ Int $\rightarrow$ SchemeVal}
\begin{lstlisting}[aboveskip=0pt]
(define nth-element
  (lambda (lst n)
    (if (null? lst)
        (report-list-too-short n)
        (if (zero? n)
            (car lst)
            (nth-element (cdr lst) (- n 1))))))
            
(define report-list-too-short
  (lambda (n)
    (eopl:error 'nth-element
                "List too short by ~s elements.~%" (+ n 1))))
\end{lstlisting}

\subsubsection{remove-first}

The procedure |remove-first| returns a given list without the first occurrence of a given symbol. If the list is empty, an empty list is returned. If the list does not contain the given symbol, the given list is returned. This functionality is available as |remv| in the Racket API.

\textbf{remove-first} : \textit{Sym $\times$ Listof(Sym) $\rightarrow$ Listof(Sym)}
\begin{lstlisting}[aboveskip=0pt]
(define remove-first
  (lambda (s los)
    (if (null? los)
        '()
        (if (eqv? (car los) s)
            (cdr los)
            (cons (car los) (remove-first s (cdr los)))))))
\end{lstlisting}

\subsubsection{occurs-free?}

The procedure |occurs-free?| evaluates a $\lambda$-calculus expression and returns |#t| if a given variable is not a bound variable in that expression. This means that it does not occur as a named parameter for any of the |lambda| expressions in the given expression.

\textbf{occurs-free?} : \textit{Sym $\times$ LcExp $\rightarrow$ Bool}
\begin{lstlisting}[aboveskip=0pt,numbers=left]
(define occurs-free?
  (lambda (var exp)
    (cond
      ((symbol? exp) (eqv? var exp))
      ((eqv? (car exp) 'lambda)
       (and
        (not (eqv? var (caadr exp)))
        (occurs-free? var (caddr exp))))
      (else
       (or
        (occurs-free? var (car exp))
        (occurs-free? var (cadr exp)))))))
\end{lstlisting}

\subsubsubsection{Lines reference}

\begin{enumerate}
\setcounter{enumi}{3}
\item If |exp| is a symbol, which is |'| followed by a letter, then the given |var| is not bound if it equals |exp|.
\item If the first element of |exp| is |'lambda|, it is a lambda expression, which we have to dissect.
\setcounter{enumi}{6}
\item |(caadr exp)| returns the parameter of the |lambda| expression, which is referred to as the first (and only) value of the head of the expression. An expression in $\lambda$-calculus can only have a single parameter for each |lambda| expression. In essence, it first performs |cdr| to remove the leading |lambda| from the expression. Then it performs |car| to obtain the head of the expression and finally it will perform another |car| to get the (first) element inside of the head.
\item |(caddr exp)| returns the body of the expression. This is only called if the previous line did not evaluate to |#f|. That means that the head of the expression did not contain |var| and then we can proceed to evaluate the body with a recursive call to |occurs-free?|.
\setcounter{enumi}{10}
\item Aside from the previous cases, $\lambda$-calculus also has expressions that come in pairs of 2. First we will extract the first value of the pair using |car| and evaluate that with a recursive call to |occurs-free?|.
\item To grab the second value of the pair, we have to first get the list of expressions symbols that does not contain the initial symbol, which we've already evaluated using |cdr| and next get the first (and only) symbol of that list using |car|. This line is only executed, if the previous line did not evaluate to |#t|.
\end{enumerate}

\subsubsection{subst}

The procedure |subst| takes 2 symbols and an s-list, which is defined as a list containing symbols and/or s-lists. The procedure should replace every occurrence of one of the given symbols in the list and all nested lists by the other given symbol, referred to as |old| and |new| respectively.

\subsubsubsection{First attempt}

\textbf{subst} : \textit{Sym $\times$ Sym $\times$ S-list $\rightarrow$ S-list}
\begin{lstlisting}[aboveskip=0pt]
(define subst
  (lambda (new old slist)
    (if (null? slist)
        '()
        (if (symbol? (car slist))
            (cons (if (eqv? (car slist) old)
                      new
                      (car slist)) (subst new old (cdr slist)))
            (cons (subst new old (car slist)) (subst new old (cdr slist)))))))
\end{lstlisting}

My first attempt can be improved by extracting the |cons| call from both expressions passed to the second |if| call, as the second expression in each |cons| call is the same.

\subsubsubsection{Textbook}

\textbf{subst} : \textit{Sym $\times$ Sym $\times$ S-exp $\rightarrow$ S-exp}
\begin{lstlisting}[aboveskip=0pt]
((define subst-in-s-exp
  (lambda (new old sexp)
    (if (symbol? sexp)
        (if (eqv? sexp old) new sexp)
        (subst new old sexp))))
\end{lstlisting}

\textbf{subst-in-s-exp} : \textit{Sym $\times$ Sym $\times$ S-list $\rightarrow$ S-list}
\begin{lstlisting}[aboveskip=0pt]
(define subst
  (lambda (new old slist)
    (if (null? slist)
        '()
        (cons
         (subst-in-s-exp new old (car slist))
         (subst new old (cdr slist))))))
\end{lstlisting}

\newpage

\subsection{Chapter 2}

\subsubsection{Environment data-structure}
\label{sec:env}

The following code represents the data-structure for an environment. |empty-env| creates a new empty environment. |extend-env| extends an environment by assigning a given value to a given variable name in a given environment. |apply-env| returns the value assigned to a given variable in a given environment.

\textbf{empty-env} : \textit{() $\rightarrow$ Env}
\begin{lstlisting}[aboveskip=0pt]
(define empty-env
  (lambda ()
    (list 'empty-env)))
\end{lstlisting}

\textbf{extend-env} : \textit{Var $\times$ SchemeVal $\times$ Env $\rightarrow$ Env}
\begin{lstlisting}[aboveskip=0pt]
(define extend-env
  (lambda (var val env)
    (list 'extend-env var val env)))
\end{lstlisting}

\textbf{apply-env} : \textit{Env $\times$ Var $\rightarrow$ SchemeVal}
\begin{lstlisting}[aboveskip=0pt]
(define apply-env
  (lambda (env search-var)
    (cond
      ((eqv? (car env) 'empty-env)
       (report-no-binding-found search-var))
      ((eqv? (car env) 'extend-env)
       (let ((saved-var (cadr env))
             (saved-val (caddr env))
             (saved-env (cadddr env)))
         (if (eqv? search-var saved-var)
             saved-val
             (apply-env saved-env search-var))))
      (else (report-invalid-env)))))

(define report-no-binding-found
  (lambda (search-var)
    (eopl:error 'apply-env "No binding for ~s" search-var)))

(define report-invalid-env
  (lambda (env)
    (eopl:error 'apply-env "Bad environment: ~s" env)))
\end{lstlisting}

\newpage

\section{Exercises}

\subsection{Chapter 1}

\subsubsection{Exercise 1.6}

If we reversed the order of checks in |nth-element|, nothing would go wrong, but we would have to perform the initial check in both branches of the |if| for the second check.

\subsubsection{Exercise 1.7}

We introduce a helper procedure |nth-element-helper| to easily preserve the original given list.

\begin{lstlisting}
(define report-list-too-short
  (lambda (lst n)
    (eopl:error 'nth-element
                "~s does not have ~s element(s).~%" l (+ (length l) (+ n 1)))))

(define nth-element-helper
  (lambda (origlst lst n)
    (if (null? lst)
        (report-list-too-short origlst n)
        (if (zero? n)
            (car lst)
            (nth-element-helper origlst (cdr lst) (- n 1))))))

(define nth-element
  (lambda (lst n)
    (nth-element-helper lst lst n)))
\end{lstlisting}

\subsubsection{Exercise 1.8}

If the last line of |remove-first| were to be replaced with |(remove-first s (cdr los))|, it would be defined as:

\textbf{remote-first} : \textit{Sym $\times$ Listof(Sym) $\rightarrow$ Listof(Sym)}\\
\textbf{usage:} |(remove-first s los)| returns a list with all elements starting after the first occurrence of the symbol |s| in the list |los|.

\subsubsection{Exercise 1.9}

We define |remove| as a procedure which removes every occurrence of a given symbol from a given list.

\begin{lstlisting}
(define remove
  (lambda (s l)
    (if (null? l)
        '()
        (if (eqv? (car l) s)
            (remove s (cdr l))
            (cons (car l) (remove s (cdr l)))))))
\end{lstlisting}

\subsubsection{Exercise 1.10}

Though typically ``or'' refers to ``inclusive or'', ``exclusive or'' also exists. The statement ``A or B'' for the inclusive variant, will evaluate to true (if it is a logical expression) when either A and/or B are true. The exclusive variant, often referred to as ``xor'', will for ``A xor B'' only evaluate to true if A or B are true, but not when both are true or false.

\subparagraph{NOTE} When we refer to bitwise operations, rather than logical ones, ``A or B'' will perform the same logic as it would for a logical expression, at the binary level. E.g. ``1001 or 0101'' evaluates to 1101, which translates to the decimal system as ``9 or 5 = 13''. ``1001 xor 0101'' evaluates to 1100, which translates to the decimal system as ``9 xor 5 = 12''.

\subsubsection{Exercise 1.11}

The |subst-in-s-exp| procedure is called with the first element of the parameter |slist| in |subst|, it only passes this back to |subst| if it is not a symbol. This means it works down into the tree of nested lists. Once the deepest nested list has been reached, recursion will automatically halt.

\subsubsection{Exercise 1.12}

The procedure |subst| defined below has |subst-in-s-exp| inlined.

\begin{lstlisting}
(define subst
  (lambda (new old slist)
    (if (null? slist)
        '()
        (cons
         (if (symbol? (car slist))
             (if (eqv? (car slist) old) new (car slist))
             (subst new old (car slist)))
         (subst new old (cdr slist))))))
\end{lstlisting}

\subsubsection{Exercise 1.13}

The procedure |map| is used to process the slist in the altered version of |subst|.

\begin{lstlisting}
(define subst
  (lambda (new old slist)
    (map (lambda (sexp)
           (if (symbol? sexp)
               (if (eqv? sexp old) new sexp)
               (subst new old sexp))) slist)))
\end{lstlisting}

\subsubsection{Exercise 1.15}

The procedure |duple| returns a list with an input value duplicated a certain number of times.

\begin{lstlisting}
(define duple
  (lambda (n x)
    (if (zero? n)
        '()
        (cons x (duple (- n 1) x)))))
\end{lstlisting}

\subsubsection{Exercise 1.16}

The procedure |invert| swaps the elements of every pair in a list of pairs.

\begin{lstlisting}
(define invert
  (lambda (lst)
    (if (null? lst)
        '()
        (cons
         (list (cadar lst) (caar lst))
         (invert (cdr lst))))))
\end{lstlisting}

\subsubsection{Exercise 1.17}

The procedure |down| wraps every top-level element of a given list in a list.

\begin{lstlisting}
(define down
  (lambda (lst)
    (if (null? lst)
        '()
        (cons
         (cons (car lst) '())
         (down (cdr lst))))))
\end{lstlisting}

\subsubsection{Exercise 1.18}

The procedure |swapper| swaps every given 2 symbols in a given list and its nested lists.

\begin{lstlisting}
(define swapper
  (lambda (s1 s2 slist)
    (map
     (lambda (elem)
       (if (symbol? elem)
           (if (eqv? elem s1)
               s2
               (if (eqv? elem s2) s1 elem))
           (swapper s1 s2 elem)))
     slist)))
\end{lstlisting}

\subsubsection{Exercise 1.19}

The procedure |list-set| replaces the element at a certain given zero-offset position in a list by a different one.

\begin{lstlisting}
(define list-set
  (lambda (lst n x)
    (if (null? lst)
        '()
        (cons
         (if (zero? n)
             x
             (car lst))
         (list-set (cdr lst) (- n 1) x)))))
\end{lstlisting}

\subsubsection{Exercise 1.20}

The procedure |count-occurrences| counts the amount of times that a given symbol is rpesent in a given list or any nested list in that list.

\begin{lstlisting}
(define count-occurrences
  (lambda (s slist)
    (if (null? slist)
        0
        (+
         (if (symbol? (car slist))
             (if (eqv? (car slist) s) 1 0)
             (count-occurrences s (car slist)))
         (count-occurrences s (cdr slist))))))
\end{lstlisting}

\subsubsection{Exercise 1.21}

The procedure |product| returns a list of pairs which represent the Cartesian product of two given lists.

\begin{lstlisting}
(define pair
  (lambda (s lst)
    (if (null? lst)
        '()
        (cons
         (list
          s
          (car lst))
         (pair s (cdr lst))))))

(define prepend
  (lambda (lst1 lst2)
    (if (null? lst1)
        lst2
        (cons
         (car lst1)
         (prepend (cdr lst1) lst2)))))

(define product
  (lambda (sos1 sos2)
    (if (or (null? sos1) (null? sos2))
        '()
        (prepend
         (pair (car sos1) sos2)
         (product (cdr sos1) sos2)))))
\end{lstlisting}

\subsubsection{Exercise 1.22}

The procedure |filter-in| returns a list containing only the top-level elements of a given list for which a given predicate evaluates to |#t|.

\begin{lstlisting}
(define filter-in
  (lambda (pred lst)
    (if (null? lst)
        '()
        (if (pred (car lst))
            (cons
             (car lst)
             (filter-in pred (cdr lst)))
            (filter-in pred (cdr lst))))))
\end{lstlisting}

\subsubsection{Exercise 1.23}

The procedure |list-index| returns the position of the first element of a given list for which a given predicate returns |#t|. If the predicate is false for every element of the list, the procedure returns |#f|.

\begin{lstlisting}
(define plus-or-false
  (lambda (a b)
    (if (or (eq? a #f) (eq? b #f))
        #f
        (+ a b))))

(define list-index
  (lambda (pred lst)
    (if (null? lst)
        #f
        (if (pred (car lst))
            0
            (plus-or-false
             (list-index pred (cdr lst))
             1)))))
\end{lstlisting}

\subsubsection{Exercise 1.24}

The procedure |every?| returns |#t| if every element of a given list evaluates to |#t| for a given predicate and |#f| if it does not.

\begin{lstlisting}
(define every?
  (lambda (pred lst)
    (if (null? lst)
        #t
        (and
         (pred (car lst))
         (every? pred (cdr lst))))))
\end{lstlisting}

\subsubsection{Exercise 1.25}

The procedure |exists?| returns |#t| if at least one element of a given list evaluates to |#t| for a given predicate and |#f| if none do.

\begin{lstlisting}
(define exists?
  (lambda (pred lst)
    (if (null? lst)
        #f
        (or
         (pred (car lst))
         (exists? pred (cdr lst))))))
\end{lstlisting}

\subsubsection{Exercise 1.26}

The procedure |up| removes the parenthesis for every top-level list in a given list.

\begin{lstlisting}
(define up
  (lambda (lst)
    (if (null? lst)
        '()
        (if (symbol? (car lst))
            (cons
             (car lst)
             (up (cdr lst)))
            (prepend
             (car lst)
             (up (cdr lst)))))))
\end{lstlisting}

|prepend| is defined in exercise 1.21.

\subsubsection{Exercise 1.27}

The procedure |flatten| removes the parenthesis for every list and nested list in a given list.

\begin{lstlisting}
(define flatten
  (lambda (lst)
    (if (null? lst)
        '()
        (if (symbol? (car lst))
            (cons
             (car lst)
             (flatten (cdr lst)))
            (prepend
             (flatten (car lst))
             (flatten (cdr lst)))))))
\end{lstlisting}

|prepend| is defined in exercise 1.21.

\subsubsection{Exercise 1.28}

The procedure |merge| takes two ordered lists of numbers and merges them into a single ordered list.

\begin{lstlisting}
(define merge
  (lambda (loi1 loi2)
    (cond
      ((and (null? loi1) (null? loi2)) '())
      ((null? loi1) (cons (car loi2) (merge loi1 (cdr loi2))))
      ((null? loi2) (cons (car loi1) (merge (cdr loi1) loi2)))
      (else
       (if (< (car loi1) (car loi2))
           (cons (car loi1) (merge (cdr loi1) loi2))
           (cons (car loi2) (merge loi1 (cdr loi2))))))))
\end{lstlisting}

\subsubsection{Exercise 1.29}

The procedure |sort| takes a list of numbers and sorts them from small to large.

\begin{lstlisting}
(define smallest
  (lambda (loi)
    (if (eq? (length loi) 1)
        (car loi)
        (if (< (car loi) (smallest (cdr loi)))
            (car loi)
            (smallest (cdr loi))))))

(define remove-first
  (lambda (s los)
    (if (null? los)
        '()
        (if (eqv? (car los) s)
            (cdr los)
            (cons (car los) (remove-first s (cdr los)))))))

(define sort
  (lambda (loi)
    (if (null? loi)
        '()
        (cons
         (smallest loi)
         (sort (remove-first (smallest loi) loi))))))
\end{lstlisting}

\subsubsection{Exercise 1.30}

The procedure |sort/predicate| takes a list of numbers and sorts them based on a given predicate.

\begin{lstlisting}
(define select
  (lambda (pred loi)
    (if (eq? (length loi) 1)
        (car loi)
        (if (pred (car loi) (select pred (cdr loi)))
            (car loi)
            (select pred (cdr loi))))))

(define sort/predicate
  (lambda (pred loi)
    (if (null? loi)
        '()
        (cons
         (select pred loi)
         (sort/predicate pred (remove-first (select pred loi) loi))))))
\end{lstlisting}

|remove-first| is defined in exercise 1.29.

\subsubsection{Exercise 1.31}

The procedure |leaf| creates the representation of a leaf in a binary tree. |interior-node| creates a named node with 2 leafs linked to it. |leaf?| returns |#t| if a given value is a leaf. |lson| returns the left leaf/subnode of a node and |rson| the right one. |contents-of| returns the value of a bintree item, which is the value of a leaf or the name of a node.

\begin{lstlisting}
(define leaf
  (lambda (v)
    v))

(define interior-node
  (lambda (s bt1 bt2)
    (list s bt1 bt2)))

(define leaf?
  (lambda (v)
   (number? v)))

(define lson
  (lambda (bintree)
    (cadr bintree)))

(define rson
  (lambda (bintree)
    (caddr bintree)))

(define contents-of
  (lambda (bintree)
    (if (leaf? bintree)
        bintree
        (car bintree))))
\end{lstlisting}

\subsubsection{Exercise 1.32}

The procedure |double-tree| doubles the value of each leaf in a tree.

\begin{lstlisting}
(define double-tree
  (lambda (bintree)
    (if (leaf? bintree)
        (leaf (* 2 bintree))
        (interior-node
         (contents-of bintree)
         (double-tree (lson bintree))
         (double-tree (rson bintree))))))
\end{lstlisting}

|leaf|, |leaf?|, |interior-node|, |lson|, |rson| and |contents-of| are defined in exercise 1.31.

\subsubsection{Exercise 1.33}

The procedure |mark-leaves-with-red-depth| replaces the value of every leaf ion a given binary tree with the number of times the symbol |'red| is found in a parent node.

\begin{lstlisting}
(define mark-leaves-depth
  (lambda (bintree s depth)
    (if (leaf? bintree)
        (leaf depth)
        (interior-node
         (contents-of bintree)
         (mark-leaves-depth
          (lson bintree)
          s
          (if (eqv? s (contents-of bintree))
              (+ 1 depth)
              depth))
         (mark-leaves-depth
          (rson bintree)
          s
          (if (eqv? s (contents-of bintree))
              (+ 1 depth)
              depth))))))

(define mark-leaves-with-red-depth
  (lambda (bintree)
    (mark-leaves-depth bintree 'red 0)))
\end{lstlisting}

|leaf|, |leaf?|, |interior-node|, |lson|, |rson| and |contents-of| are defined in exercise 1.31.

\subsubsection{Exercise 1.34}

The procedure |path| returns a list of |'left| and |'right| symbols to indicate a path to where a given value can be found in a binary search tree.

\begin{lstlisting}
(define contains
  (lambda (n bst)
    (if (null? bst)
        #f
        (or
         (eq? n (car bst))
         (or
          (contains n (cadr bst))
          (contains n (caddr bst)))))))

(define path
  (lambda (n bst)
    (if (or (null? bst) (eq? n (car bst)))
        '()
        (if (contains n (cadr bst))
            (cons
             'left
             (path n (cadr bst)))
             (if (contains n (caddr bst))
                 (cons
                  'right
                  (path n (caddr bst)))
                 '())))))
\end{lstlisting}

\subsubsection{Exercise 1.35}

The procedure |number-leaves| takes a given bintree and returns a tree with the same structure where the leaves have been replaced with their count from 0 to n-1, starting from the left-most leaf going to the right-most leaf in the tree.

\begin{lstlisting}
(define count-leaves
 (lambda (bintree)
   (if (leaf? bintree)
       1
       (+
        (count-leaves (lson bintree))
        (count-leaves (rson bintree))))))

(define number-bintree
  (lambda (bintree n)
    (if (leaf? bintree)
        (leaf n)
        (interior-node
         (contents-of bintree)
         (number-bintree (lson bintree) n)
         (number-bintree (rson bintree) (+ n (count-leaves (lson bintree))))))))

(define number-leaves
  (lambda (bintree)
    (number-bintree bintree 0)))
\end{lstlisting}

|leaf|, |interior-node|, |leaf?|, |lson|, |rson| and |contents-of| are defined in exercise 1.31.

\subsubsection{Exercise 1.36}

The procedure |number-elements| takes a list of elements and returns a list of pairs, where every pair is the index of an element from the original list, paired with the element itself.

\begin{lstlisting}
(define f
  (lambda (lst rlst n)
    (cons
     (list
      (+ (car lst) n)
      (cadr lst))
     (if (null? rlst)
         '()
         (f
          (car rlst)
          (cdr rlst)
          1)))))

(define g
  (lambda (lst rlst)
    (f lst rlst 0)))

(define number-elements
  (lambda (lst)
    (if (null? lst) '()
        (g (list 0 (car lst)) (number-elements (cdr lst))))))
\end{lstlisting}

\newpage

\subsection{Chapter 2}

\subsubsection{Exercise 2.1}

The procedure |factorial| takes a base and a number is bigits for that base, for which it will determine the factorial for that number in bigits for that same base. The base operations for the bigits datatype are defined as |zero|, |is-zero?|, |successor| and |predecessor|. To build |factorial|, the mathematical operations |plus| and |times| have also been defined.

As the base gets larger, the time required to compute the |factorial| decreases. This can be explained by the fact that less bigits will be required to represent larger numbers, which means that there have to be less recursive |successor| and |predecessor| calls made to calculate the result.

\begin{lstlisting}
(define zero (lambda () '()))

(define is-zero? (lambda (bigits) (null? bigits)))

(define successor
  (lambda (N bigits)
    (if (is-zero? bigits)
        '(1)
        (if (eq? (car bigits) (- N 1))
            (cons 0 (successor N (cdr bigits)))
            (cons (+ 1 (car bigits)) (cdr bigits))))))

(define predecessor
  (lambda (N bigits)
    (if (eq? (car bigits) 0)
        (cons (- N 1) (predecessor N (cdr bigits)))
        (if (and
             (eq? (- (car bigits) 1) 0)
             (null? (cdr bigits)))
            (zero)
            (cons (- (car bigits) 1) (cdr bigits))))))

(define plus
  (lambda (N bigits1 bigits2)
    (if (is-zero? bigits2)
        bigits1
        (plus N (successor N bigits1) (predecessor N bigits2)))))

(define times
  (lambda (N bigits1 bigits2)
    (if (is-zero? bigits2)
        (zero)
        (plus N bigits1 (times N bigits1 (predecessor N bigits2))))))

(define factorial
  (lambda (N bigits)
    (if (is-zero? bigits)
        '(1)
        (times N bigits (factorial N (predecessor N bigits))))))
\end{lstlisting}

\subsubsection{Exercise 2.4}

The procedure |empty-stack| returns an empty stack. |push| adds a given element to the top of a given stack. |pop| removes an element from the top of a given stack. |top| returns the element at the top of a given stack. |empty-stack?| returns |#t| if a stack is empty.

|empty-stack|, |push| an |pop| are constructors, as they create or modify a given stack. |top| and |empty-stack?| are observers as they only retrieve data from a given stack.

\begin{lstlisting}
(define empty-stack (lambda () '()))

(define push
  (lambda (val stack)
    (if (list? stack)
        (cons val stack)
        (eopl:error 'push "Invalid stack"))))

(define pop
  (lambda (stack)
    (cond
      ((empty-stack? stack) (eopl:error 'pop "Stack is empty"))
      ((list? stack) (cdr stack))
      (else (eopl:error 'pop "Invalid stack")))))

(define top
  (lambda (stack)
    (cond
      ((empty-stack? stack) (eopl:error 'top "Stack is empty"))
      ((list? stack) (car stack))
      (else (eopl:error 'top "Invalid stack")))))
    
(define empty-stack?
  (lambda (stack)
    (null? stack)))
\end{lstlisting}

\subsubsection{Exercise 2.5}

The procedure |empty-env|, |extend-env| and |apply-env| procedures represent an environment as an association list (a-list), with the empty environment as an empty list.

\begin{lstlisting}
(define empty-env (lambda () '()))

(define extend-env
  (lambda (var val env)
    (list (list var val) env)))

(define apply-env
  (lambda (env search-var)
    (cond
      ((null? env) (report-no-binding-found search-var))
      ((and (and (pair? env) (pair? (car env))) (list? (cadr env)))
       (if (eqv? (caar env) search-var)
           (cadar env)
           (apply-env (cadr env) search-var)))
      (else (report-invalid-env env)))))
\end{lstlisting}

|report-no-binding-found| and |report-invalid-env| are defined in the code for environments in section \ref{sec:env}.

\subsubsection{Exercise 2.6}

A first alternate representation is a linked list, as seen in exercise 2.5. An alternative to the a-list is a list of pairs:

\begin{lstlisting}
(define empty-env (lambda () '()))

(define extend-env
  (lambda (var val env)
    (cons (list var val) env)))

(define apply-env
  (lambda (env search-var)
    (cond
      ((null? env) (report-no-binding-found search-var))
      ((and (and (list? env) (pair? (car env))) (list? (cdr env)))
       (if (eqv? (caar env) search-var)
           (cadar env)
           (apply-env (cdr env) search-var)))
      (else (report-invalid-env env)))))
\end{lstlisting}

Another alternative is a sequential listing of variables and values, which are not specifically grouped, just listed in sequence:

\begin{lstlisting}
(define empty-env (lambda () '()))

(define extend-env
  (lambda (var val env)
    (cons var (cons val env))))

(define apply-env
  (lambda (env search-var)
    (cond
      ((null? env) (report-no-binding-found search-var))
      ((and (and (list? env) (symbol? (car env))) (not (null? (cdr env))))
       (if (eqv? (car env) search-var)
           (cadr env)
           (apply-env (cddr env) search-var)))
      (else (report-invalid-env env)))))
\end{lstlisting}

|report-no-binding-found| and |report-invalid-env| are defined in the code for environments in section \ref{sec:env}.

\subsubsection{Exercise 2.8}

The procedure |empty-env?| returns |#t| if a given environment is empty. This procedure works for the original interface defined in section \ref{sec:env}, as well as those defined in exercises 2.5 and 2.6.

\begin{lstlisting}
(define empty-env?
  (lambda (env)
    (equal? env (empty-env))))
\end{lstlisting}

|empty-env| is defined in the sections and exercises mentioned above the code.

\subsubsection{Exercise 2.9}

The procedure |has-binding?| returns |#t| if a given environment contains a binding for a given variable. The following implementation works with the interface defined in section \ref{sec:env}.

\begin{lstlisting}
(define has-binding?
  (lambda (env search-var)
    (cond
      ((eqv? (car env) 'empty-env) #f)
      ((eqv? (car env) 'extend-env)
       (let ((saved-var (cadr env))
             (saved-val (caddr env))
             (saved-env (cadddr env)))
         (or (eqv? search-var saved-var)
             (has-binding? saved-env search-var))))
      (else (report-invalid-env)))))
\end{lstlisting}

An alternative implementation for a-lists used in exercise 2.5 can also be defined.

\begin{lstlisting}
(define has-binding?
  (lambda (env search-var)
    (cond
      ((null? env) #f)
      ((and (and (list? env) (pair? (car env))) (list? (cdr env)))
       (or (eqv? (caar env) search-var)
           (has-binding? (cdr env) search-var)))
      (else (report-invalid-env env)))))
\end{lstlisting}

|report-invalid-env| is defined in the code in section \ref{sec:env}.

\subsubsection{Exercise 2.10}

The procedure |extend-env*| extends a given environment with a given list of variables and a given list of their respective values. The procedure defined below is intended for use with the a-list interface defined in exercise 2.5.

\begin{lstlisting}
(define extend-env*
  (lambda (vars vals env)
    (if (or (null? vars) (null? vals))
            env
            (list
             (list
              (car vars)
              (car vals))
             (extend-env* (cdr vars) (cdr vals) env)))))
\end{lstlisting}

\subsubsection{Exercise 2.11}

The procedure |extend-env*| extends a given environment with a given list of variables and a given list of their respective values. This particular implementation uses a backbone structure to perform the operation of extending the environment in constant time. The rest of the interface for the environment based on a backbone structure has also been provided as |empty-env|, |extend-env| and |apply-env|.

\begin{lstlisting}
(define empty-env (lambda () '()))

(define extend-env
  (lambda (var val env)
    (list (list (list var) (list val)) env)))

(define find-env-value
  (lambda (vars vals search-var)
    (if (or (null? vars) (null? vals))
        #f
        (if (eqv? (car vars) search-var)
            (car vals)
            (find-env-value (cdr vars) (cdr vals) search-var)))))

(define apply-env
  (lambda (env search-var)
    (cond
      ((null? env) (report-no-binding-found search-var))
      ((and (and (and (and (pair? env) (pair? (car env)))(list? (caar env)))
                 (list? (cadar env))) (list? (cadr env)))
       (let ([val (find-env-value (caar env) (cadar env) search-var)])
         (if (eq? val #f)
             (apply-env (cadr env) search-var)
             val)))
      (else (report-invalid-env env)))))
      
(define extend-env*
  (lambda (vars vals env)
    (if (or (null? vars) (null? vals))
            env
            (list (list vars vals) env))))
\end{lstlisting}

|report-no-binding-found| and |report-invalid-env| are defined in the code for environments in section \ref{sec:env}.

\newpage

\section{Racket (Scheme) Language Reference}

\subsection{API}

\subsubsection{car}

\begin{lstlisting}[frame=none]
(car l)
\end{lstlisting}

\textbf{Description}

Returns the first element of a given list |l|.

\textbf{Examples}

\begin{lstlisting}
> (car '(a b (c d) e))
a
\end{lstlisting}

\subsubsection{cdr}

\begin{lstlisting}[frame=none]
(cdr l)
\end{lstlisting}

\textbf{Description}

Returns the n-1 last elements of a given list |l| as a new list, where n is the length of the list |l|.

\textbf{Examples}

\begin{lstlisting}
> (cdr '(a b (c d) e))
(b (c d) e)
\end{lstlisting}

\subsubsection{cons}

\begin{lstlisting}[frame=none]
(cons a b)
\end{lstlisting}

\textbf{Description}

Returns a pair containing |a| and |b| if they are both symbols. If |b| is a list, it will return that list with |a| appended at the start.

\textbf{Examples}

\begin{lstlisting}
> (cons 'a 'b)
'(a . b)
\end{lstlisting}

\begin{lstlisting}
> (cons 'a '(b))
'(a b)
\end{lstlisting}

\subsubsection{define}

\begin{lstlisting}[frame=none]
(define id expr)
(define (head args) body)
\end{lstlisting}

\textbf{Description}

The first definition binds |id| to the result of |expr|. The second one binds a |head| with |args| to a |body|.

\textbf{Examples}

\begin{lstlisting}
(define x 10)

> x
10
\end{lstlisting}

\begin{lstlisting}
(define (f x)
  (+ x 1))
  
> (f 10)
11
\end{lstlisting}


\subsubsection{eq?}

\begin{lstlisting}[frame=none]
(eq? a b)
\end{lstlisting}

\textbf{Description}

Returns |#t| if 2 given values |a| and |b| refer to the same object. Basic values (and strings) will be compared by value, as they are not represented by objects. To compare by value, refer to |eqv?|.

\textbf{Examples}

\begin{lstlisting}
> (eq? 1 1)
#t
\end{lstlisting}

\begin{lstlisting}
> (eq? 'a 'a)
#f
\end{lstlisting}

\begin{lstlisting}
> (eq? '(a b) '(a b))
#f
\end{lstlisting}

\begin{lstlisting}
> (eq? "abc" "abc")
#t
\end{lstlisting}


\subsubsection{eqv?}

\begin{lstlisting}[frame=none]
(eqv? a b)
\end{lstlisting}

\textbf{Description}

Returns |#t| if 2 given values |a| and |b| are equivalent and |#f| if they are not. This means that the values are equal, but is not to be confused with the procedure |eq?|. This can be used as an alternative to |=| to compare any type of values, as |=| only compares numbers. |eqv?| can not be used to verify if 2 lists are equal, for this refer to |equal?|.

\textbf{Examples}

\begin{lstlisting}
> (eqv? 1 1)
#t
\end{lstlisting}

\begin{lstlisting}
> (eqv? 'a 'a)
#t
\end{lstlisting}

\begin{lstlisting}
> (eqv? '(a b) '(a b))
#f
\end{lstlisting}

\begin{lstlisting}
> (eqv? "abc" "abc")
#t
\end{lstlisting}

\subsubsection{if}

\begin{lstlisting}[frame=none]
(if test-expr then-expr else-expr)
\end{lstlisting}

\textbf{Description}

Evaluates |test-expr| and evaluates |then-expr| if the result was anything but |#f|, otherwise |else-expr| is evaluated.

\textbf{Examples}

\begin{lstlisting}
> (if (= 1 2) 3 4)
4
\end{lstlisting}

\subsubsection{lambda}

\begin{lstlisting}[frame=none,belowskip=0pt]
(lambda args body)
\end{lstlisting}
\begin{lstlisting}[frame=none,language=SchemeLambda,aboveskip=0pt]
(lambda args body)
\end{lstlisting}

\textbf{Description}

Produces an anonymous procedure with given |args| and |body|.

\textbf{Examples}

\begin{lstlisting}
> ((lambda (x) x) 10)
10
\end{lstlisting}

\subsubsection{let}

\begin{lstlisting}[frame=none]
(let ([id val-expr] ...) body)
\end{lstlisting}

\textbf{Description}

Assigns |val-expr| to |id| for a set of these pairs. After the assignment, |body| is evaluated with these new identifiers.

\textbf{Examples}

\begin{lstlisting}
> (let ([x 10] [y 5]) (+ x y))
15
\end{lstlisting}

\subsubsection{list-ref}

\begin{lstlisting}[frame=none]
(list-ref l p)
\end{lstlisting}

\textbf{Description}

Returns the element of a given list |l| at position |p|. The position if offset at 0.

\textbf{Examples}

\begin{lstlisting}
> (list-ref '(a b (c d) e) 2)
(c d)
\end{lstlisting}

\subsubsection{null?}

\begin{lstlisting}[frame=none]
(null? l)
\end{lstlisting}

\textbf{Description}

Returns |#t| if |l| is an empty list |'()| and |#f| is it is not.

\textbf{Examples}

\begin{lstlisting}
> (null? '())
#t
\end{lstlisting}

\subsubsection{symbol?}

\begin{lstlisting}[frame=none]
(symbol? s)
\end{lstlisting}

\textbf{Description}

Returns |#t| if |l| is a symbol and |#f| is it is not.

\textbf{Examples}

\begin{lstlisting}
> (symbol? 'a)
#t
\end{lstlisting}

\begin{lstlisting}
> (symbol? 1)
#f
\end{lstlisting}

\subsubsection{zero?}

\begin{lstlisting}[frame=none]
(zero? z)
\end{lstlisting}

\textbf{Description}

Returns |#t| if |z| is zero and |#f| is it is not.

\textbf{Examples}

\begin{lstlisting}
> (zero? 0)
#t
\end{lstlisting}

\end{document}

